% The captions are included in generate_plots.tex so that the size of the figures can be adjusted.
% However, this means that the same caption would have to be maintained both there and where the figure is included in the main text.
% The solution is to break them out into this file.

\newcommand{\figellipticcaption}{%
    \caption{The integrals $E_1$, $E_2$ and $E_3$ (left to right), shown as functions of the nome $q$ over the unit disk in the complex plane.
    For each function, its value at $t=0$ (limit as $q\to1$) has been subtracted to make the features more visible.
    The hue represents the complex argument (with \textbf{\color{red}red} and \textbf{\color{cyan}cyan} for positive and negative real values, respectively), while the shaded contours represent the magnitude logarithmically: with each successive contour, $|t|$ gets a factor of $2$ larger (brighter side) or smaller (darker side).
    Each $t\neq 0$ corresponds to infinitely many $q$ via \cref{eq:q-to-t}, but the wedge-shaped region outlined in white corresponds to a single copy of the complex $t$-plane, with $t=0$ at its rightmost extremity; this is explored further in \cref{sec:t-tau-q}.
    Note the branch cut on $q\in[0,1]$, which corresponds to two copies of $t\in[16,\infty)$ (see \cref{fig:q}).
    }}


\newcommand{\figtaucaption}{%
    \caption{\textbf{Left}: the mapping $\tau\mapsto t$ visualized in the complex plane, covering one period in the real direction.
%     The hue represents the complex argument (with \textbf{\color{red}red} and \textbf{\color{cyan}cyan} for positive and negative real values, respectively), while the shaded contours represent the magnitude logarithmically: with each successive contour, $|t|$ gets a factor of $2$ larger (brighter side) or smaller (darker side).
    The complex values are represented as in \cref{fig:elliptic}.
    The $\times$-shaped contour intersections are saddle points at $t=4$ (for instance at \mbox{$\tau=\pm\frac14+\frac{i\sqrt3}{12}$}) and $t=16$ (for instance at \mbox{$\tau=\pm\frac12+\frac{i\sqrt3}{6}$}).
    \textbf{Right}: a sketch (to scale) of the $\Re(\tau)>0$ half of the same plot, showing a selection of the infinitely many $\tau$-plane images of the real $t$-axis, colored according to which kinematic region they reproduce: spacelike subthreshold (\mbox{$t<0$}, \textbf{\color{spacelike}blue}), timelike subthreshold (\mbox{$0<t<4$}, \textbf{\color{subthr}orange}), two-pion (\mbox{$4<t<16$}, \textbf{\color{twopi}green}) and multi-pion (\mbox{$16<t$}, \textbf{\color{fourpi}violet}).
    The shaded area (if extended upward to infinity) contains all $t$ with $\Im(t)>0$ exactly once, with contours indicating (bottom to top) $\Im(t)=0.001$, $0.1$, $0.5$, and $1$; its boundary (bold) contains all real $t$ exactly once, ordered counterlockwise.}}

\newcommand{\figqcaption}{%
    \caption{Like \cref{fig:tau}, but showing $q\mapsto t$ over the unit disk.
            The lines in the sketch have been extended onto the $\Re(\tau)<0$ part (lower half of the disk), and thereby produce the outline used in \cref{fig:elliptic}.
            Note how inside the shaded region, $q$ approaches $1$ extremely slowly as $|t|$ approaches $0$, ensuring good convergence of our sums almost everywhere.
        }}

\newcommand{\figtsuncaption}{%
    \caption[]{Like \cref{fig:tau}, but showing {\tikzexternaldisable $t_\sunset\mapsto t$} as defined in \cref{eq:t-to-tau}.
        In the sketch on the right, the shaded area is $\Omega_>$, and the two inverse images of the real $t$-line are shown with bold colored lines as in \cref{fig:tau}.
        The dotted rectangle indicates the region shown in \cref{fig:rho24} (bottom row) and \cref{fig:rhoF}.
        The thin gray lines are the same contours of constant $\Im(t)$ as in \cref{fig:tau}, and the dashed ones indicate their images under the involutory mapping $z\mapsto (9z-9)/(z-9)$, which is how they appear in $\varpi_r$.}}


\newcommand{\figrhotwofourcaption}{%
    \caption{
        The construction of $\rho_2(z)$ (top row) and $\rho_4(z)$ (bottom row).
        \textbf{Left}: The principal branch of $\sqrt{(z-4)(z-16)}$ and $\sqrt[4]{z^3 - 9z^2 + 3z - 3}$, respectively, visualized in the complex plane using colors and contours like in \cref{fig:tau}.
        \textbf{Middle}: The corrected functions $\rho_2(z)$ and $\rho_4(z)$, respectively.
        \textbf{Right}: A schematic picture of their relation.
        Colored lines indicate where the argument of the root [$(z-4)(z-16)$ and $z^3 - 9z^2 + 3z - 3$, respectively] is real: \textbf{\color{plotV} red} for positive values and \textbf{\color{plotII} cyan} for negative.
        Dotted lines do not cause branch cuts.
        Solid lines cause cuts that are removed by analytically continuing to different branches, starting on the principal branch in the region marked by $\star$ and picking up roots of unity ($-1$ and $+i$, respectively) when crossing a line in the direction indicated by the arrows.
        Zigzag lines are left as branch cuts also in $\rho_2(z)$ or $\rho_4(z)$.
        The dotted circle indicates the boundary between $\Omega_<$ and $\Omega_>$ (compare \cref{fig:tsun}).
%             It can be noted that a version with cuts only along the dotted red lines is possible by assigning sheets differently, but that is not suitable for $t\mapsto\tau$.
        }}

\newcommand{\figrhoFcaption}{%
    \caption[The function $\rho_F(z)-F_1(1)$, visualized over the same region as $\rho_4(z)$ in \cref{fig:rho24}.]{
        \textbf{Top}:
        The function $\rho_F(z)-F_1(1)$, visualized over the same region as $\rho_4(z)$ in \cref{fig:rho24}.
        Subtracting the constant $F_1(1)$ makes the features more visible.
        The upper half-plane shows $\rho_F(z)$ as used for the
        % $\map{t}{\tau}$
        $t\mapsto\tau$
        relation, and the lower shows the uncorrected $F_1\big(1728/j_\sunset(z)\big)$ using the principal branch of the hypergeometric function.
        Both are symmetric (up to phases) across the real axis.
        The insets show the region around $z=9$, which for the uncorrected function contains the only discontinuity outside $\Omega_<$, making its removal the most important part of the construction.
%         The discontinuity in the magnified region around $z=9$ has a particularly devastating effect on
%         $\map{t}{\tau}$ since it is contained in $\Omega_>$, not $\Omega_<$,
%         % $t\mapsto\tau$,
%         % since the image of the real line under the mapping
%         % $\map{t}{t_\sunset}$
%         % $t\mapsto t_\sunset$
%         % passes through it in a way that
%         and unless remedied, it
%         leads to very problematic interplay between $\varpi_c$ and $\varpi_r$.
        \textbf{Bottom}:
        A schematic picture of the construction of the correct $\rho_F(z)$,
        using the same line conventions as in \cref{fig:rho24}: dotted for no cut, solid for removed cut, zigzag for remaining cut.
        Negative real $z$ values only produce a cut for $F_2(z)$ and therefore do not affect the principal branch ($\star$), which is purely $F_1(z)$.
        Positive real $z$ values only produce a cut when \mbox{$|1728/j_\sunset(z)|\geq 1$}, which holds within the \textbf{\color{realextra} orange} area.
        The monodromies described in the text are used when crossing the solid lines to keep $\rho_F(z)$ continuous: $C_0$ ($\rightarrow$) or $C_1$ ($\leftrightarrow$) is applied when going in the direction of the arrows (recall that $C_1$ is involutory, hence the bidirectional arrow).
        }}

\newcommand{\figrhoFconstrcaption}{%
    \caption[ A schematic version of \cref{fig:rhoF}, showing the construction of the correct $\rho_F(z)$.]{
        A schematic version of \cref{fig:rhoF}, showing the construction of the correct $\rho_F(z)$.
        The lines indicate where $1728/j_\sunset(z)$ (the argument of the hypergeometric functions) is real, using the same conventions as in \cref{fig:rho24}: dotted for no cut, solid for removed cut, zigzag for remaining cut.
        Negative real values only produce a cut for $F_2(z)$ and therefore do not affect the principal branch ($\star$), which is purely $F_1(z)$.
        Positive real values only produce a cut when $|1728/j_\sunset(z)|\geq 1$, which holds within the shaded orange area.
        The monodromies described in the text are used when crossing the solid lines to keep $\rho_F(z)$ continuous: $C_0$ ($\rightarrow$) or $C_1$ ($\leftrightarrow$) is applied when going in the direction of the arrows (recall that $C_1$ is its own inverse, hence the bidirectional arrow).
        }}

\newcommand{\fignumericexamplecaption}{%
    \caption{
        Example of the process of computing $\bar E_5$ and $\bar E_6$ at $t_\ex=4.2+i\epsilon$ following \cref{eq:Ebar-detailed}, with the final summation following \cref{eq:Ebar-schematic}.
        The integration contours are shown in the $\tau$ and $\xi$ planes, which are represented as in \cref{fig:tau}.
        Filled triangles indicate a complex integration contour, and empty ones a real integral toward infinity (using $\theta$ for that along the imaginary $\tau$ axis).
        The shaded area indicates the ``too close to the cut'' zone described in the definition of $\mathcal C_\xi(\beta)$.
        For integrals using \cref{eq:hybrid-integral}, the parts from below ($<$) and above ($>$) the cutoff $\chi$ are indicated.
        % All numbers are truncated; see \cref{tab:numeric-example} for the full numbers.
        All numbers are much more precise than shown here, with the largest errors [those of $\IE{n}(\beta_\star)$] being on the order of $10^{-6}$.
        % Note the presence of large cancellations, with several intermediate quantities being several orders of magnitude larger than the end result; this is even more striking when $t\in[0,4)$, where the large imaginary parts of the contours cancel to give a real result.
        % (Of course, if the equivalent method described at the end of \cref{sec:spacelike} is used, everything is manifestly real).
        }}

\newcommand{\fighybridcaption}{%
    \caption{
        The four usages of \cref{eq:hybrid}, evaluated at $\beta=\beta_\star$ and shown as functions of the cutoff $\chi$.
%           In each case, the top panel shows the total value (\linesample{totval}) along with that of the ``head'' (\linesample{headval}), i.e., $\xi<\chi$ numerically integrated, and the ``tail'' (\linesample{tailval}), i.e., $\xi>\chi$ using a series expansion.
        The bottom panel shows the absolute error estimate of each part.
        For easier comparison, all values (but not the errors) have been divided by the corresponding total value of the function evaluating using our preferred $\chi$.
        This $\chi$, which is what we use elsewhere, is marked with a dot.
        }}

\newcommand{\figmonodromycaption}{%
    \caption{The contours used in the monodromy construction, with cuts drawn as zigzag lines for the different functions involved.
        \textbf{Left}: $F_1(z)$, cut on $[1,\infty)$.
        \textbf{Middle}: $F_2(z)$, cut on $(-\infty,0]$.
        \textbf{Right}: $\hypergeometric abc{\frac1z}$, showing how the combination of $C_0$ and $C_1$ results in a contour that encircles the cut without crossing it.
        }}

\newcommand{\figJbubcaption}{%
    \caption[The function Jbub.]{
        The functions $\Jbub{n}(t)$, normalized so that the magnitude is roughly constant in $n$.
        Real parts are shown as solid lines, and imaginary parts as dashed lines; colors indicate different $n$ as per the plot legend.
        Below the 2-pion threshold ($t=4$), the imaginary parts are identically zero and are not shown.
        Above threshold, $t$ has been given a small positive imaginary part.
        }}

\newcommand{\figmasterscaption}{%
    \caption[The six elliptic master integrals plotted as functions of $t$.]{
        The six elliptic master integrals plotted as functions of $t$, showing the real (\Resample) and imaginary (\Imsample) parts, as well as the fifth-order series expansion around $t=0$ (\seriessample).
%         The two- and four-pion thresholds are drawn as vertical dotted lines.
%         Below each plot is shown the relative difference between \texttt{pySecDec} (lines with error bands) and our implementation, whose error is not visible at the scale of the plots.
%         Also shown is the relative difference between our implementation and the series expansion.
%         Reinterpreting the below-threshold imaginary part as errors (like in \cref{fig:abserr}) gives error bands much taller than the plots.
        At the scale of the plots, the difference between our implementation and \texttt{pySecDec} is not visible (but see \cref{fig:reldiff}), nor are the error bands.
        However, the absolute uncertainty of each line is shown on a logarithmic scale below each plot, using \texttt{pySecDec}'s uncertainty; the uncertainty using our method is shown as a shaded area.
        For the series expansion, the truncation error is estimated as the magnitude of the last non-truncated term.
        }}

\newcommand{\figreldiffcaption}{%
    \caption[Symmetric log plots of the relative difference between the three different implementations of the master integrals]{
        Symmetric log plots of the relative difference between our results and those of \texttt{pySecDec} and the series expansion, shown similarly to \cref{fig:masters}.
        By relative difference, we mean $[x - x_\text{ref}]/x_\text{ref}$ where $x$ is some quantity and $x_\text{ref}$ is the same quantity as computed using our method.
        Below threshold, the imaginary part is known to be zero and the relative difference makes little sense, so in the parts of the plots that are shaded ({\setbox0=\hbox{Xg}\tikzset{external/export next=false}\tikz[baseline=\dp0]{\draw[gray, pattern=north east lines, pattern color=gray] (0,0) rectangle (\dimexpr\ht0+\dp0\relax,\dimexpr\ht0+\dp0\relax);}}) we display the actual quantities $x$ rather than the relative difference.
        Our values (whose relative difference is zero by construction) and their error bands are shown in gray.
        }}

\newcommand{\figtimecaption}{%
    \caption[Time consumption, measured in seconds per master integral evaluation, needed by the different evaluation methods for different $t$.]{
        Time, measured in seconds per sample, needed to obtain the data behind \cref{fig:masters,fig:reldiff}.
        Solid lines indicate our implementation, and dashed lines indicate \texttt{pySecDec}; colors represent different masters as per the plot legend.
        Note that our implementation is single-threaded, whereas \texttt{pySecDec} made full use of the 16-core processor on which this data was collected.
        }}

\newcommand{\figPiEcaption}{%
    \caption[$\bar\Pi_E$]{
        The real (solid) and imaginary (dashed) parts of $\bar\Pi_E(t)$, and its fifth-order series expansion around $t=0$ (dotted), which has radius of convergence $4$.
        The error bands are not visible on the scale of the plot, but the absolute uncertainties are shown on a logarithmic scale below the main panel.
        Also shown are the contributions from individual master integrals (i.e., the term in \cref{eq:PiE} consisting of that integral times a Laurent polynomial of $t$), as well as the $E_i$-independent remainder.
        The imaginary parts are only shown above their respective thresholds.
        The error bands are too small to be seen at the scale of the plot, but the absolute uncertainties are shown on a logarithmic scale below the main plot.
        }}

\newcommand{\figPiJcaption}{%
    \caption[$\bar\Pi_J$]{
        $\bar\Pi_J(t)$, shown similarly to $\bar\Pi_E(t)$ in \cref{fig:PiE}.
        Here, the contributions are separated based on proportionality to different products of $\Jbub{n}$.
        Also shown is the purely real $\bar\Pi_\zeta(t)$ (dotted).
        The uncertainties are minuscule and are not shown.
        }}


\newcommand{\figtauerrcaption}{%
    \caption[]{
        The numerical error $\delta\tau\coloneq\map{t}{\tau}\big[\map{\tau}{t}(\tau)\big]-\tau$ obtained when mapping $\tau$ to $t$ and back, shown in the $\tau$ plane similarly to \cref{fig:tau}.
        Since $\map{\tau}{t}$ is exact to within negligible arithmetic noise
        % (128-bit floating point numbers are used here, or $\gtrsim35$ digits of precision),
        this effectively shows the precision of $\map{t}{\tau}$, which is clearly dominated by a smooth systematic error rather than noise.
        The left panel illustrates the phase of $\delta\tau$ as in figs.~\ref{fig:tau},~\ref{fig:q}, etc., whereas the right shows only the absolute value to more clearly demonstrate its size.
        The jump in error marked by the arrow corresponds to the argument of $\varpi_r$ crossing the branch cut near $z=9$ that is shown magnified in \cref{fig:rhoF}; although the discontinuity of $\rho_F$ has been removed using the monodromy, the introduction of $F_2$ of course changes its uncertainty.
        Note also the conspicuous row of four zeroes (dark spots in the right panel).
        These lie along the curve $|t_\sunset-9|=\sqrt{72}$, on which \cref{eq:t-to-tau} reduces to $i\varpi_c(t_\sunset)/\varpi_r(t_\sunset^*)$; evidently, this allows the errors to more or less cancel.}}

